\hypertarget{abs-day-monday}{}
\daypalettebanner{Monday}{mondaybg}{mondayaccent}{June 29, 2026}

\section*{Theme 1: Black Hole Demographics and Scaling Relations}
\addcontentsline{toc}{section}{Theme 1: Black Hole Demographics and Scaling Relations}

\TopicNeedSpace{10\baselineskip}
\hypertarget{abs-8}{}
\subsubsection*{Time Domain Studies of Supermassive Black Holes}
\addcontentsline{toc}{subsubsection}{Time Domain Studies of Supermassive Black Holes}
\textbf{Speaker:} Erin Kara \invitedbadge{}\\

The growth of supermassive black holes via accretion is inherently a turbulent process, causing variability in the observed multi-wavelength emission. We can use this variability to infer properties of the accretion flow (e.g. the disk geometry, accretion rate, density) and the black hole itself (e.g. its mass and spin). In this talk, I will review AGN variability studies, including recent reverberation mapping results, and what can be learned from more exotic transient accretion episodes, such as Tidal Disruption Events and Quasi Periodic Eruptions.

\TopicNeedSpace{10\baselineskip}
\hypertarget{abs-9}{}
\subsubsection*{Spatially Resolving the Nuclear Regions of Quasars with GRAVITY+}
\addcontentsline{toc}{subsubsection}{Spatially Resolving the Nuclear Regions of Quasars with GRAVITY+}
\textbf{Speaker:} Taro Shimizu \invitedbadge{}\\

Interferometry with GRAVITY at the VLTI has opened a new window into studies of the nuclear regions around accreting supermassive black holes. In this talk, I will review the work of our group, focussing on our resolved observations of the broad line region (BLR) in both local and high redshift AGN and quasars. Through kinematic modelling of GRAVITY data we have independently measured the BLR size and SMBH mass for 7 local AGN and 7 high redshift quasars (z > 2). We find deviations of the BLR size luminosity relation at both high luminosity and high Eddington ratio as well as a high rate of BLRs dominated by outflow motions. In addition, nearly all of our high redshift quasars are accreting at super-Eddington levels. Combined with follow up NOEMA and ALMA data, we determine host galaxy masses that indicate our high redshift quasars are overmassive, similar to the recently discovered AGN by JWST. With the new LGS systems of GRAVITY+, we further have recently resolved the nuclear [OIII] emitting region with evidence for both a nuclear rotating disk and large-scale outflow. Finally, I will end my talk with an outlook to the future of interferometry with kilometer scale baselines.

\TopicNeedSpace{10\baselineskip}
\hypertarget{abs-14}{}
\subsubsection*{Advancing Black Hole Mass Measurements in Nearby Massive Galaxies with JWST and Triaxial Stellar-dynamical Models}
\addcontentsline{toc}{subsubsection}{Advancing Black Hole Mass Measurements in Nearby Massive Galaxies with JWST and Triaxial Stellar-dynamical Models}
\textbf{Speaker:} Refa Al-Amri \\
\textbf{Affiliation:} Texas A\&M University \textperiodcentered{} Student (BSc/MSc/PhD)\\
About 150 supermassive black hole (BH) masses have been dynamically measured in nearby galaxies, but the local BH census remains highly incomplete. As a result, the slope, scatter, and even the shape of the BH-host galaxy relations are still uncertain, particularly at the high-mass end. In addition, only a small number of studies have directly compared BH masses derived from different dynamical tracers or modeling assumptions, with mixed results. In this talk, we focus on two local ellipticals, M87 and NGC 4261, which anchor the sparsely populated upper end of the BH mass distribution. A wide range of BH mass measurements has been reported for M87 over the past four decades using multiple methods, whereas NGC 4261 was the target of one of the earliest HST ionized-gas dynamical studies and now has a very precise BH mass from ALMA molecular-gas observations. We present exquisite high angular resolution spectra at the center of both galaxies from the JWST NIRSpec IFU, large-scale stellar kinematics measured with Keck KCWI and VLT MUSE, and stellar-dynamical models that make use of a recently revamped triaxial orbit code and modern search strategies to simultaneously determine mass and galaxy intrinsic shape parameters. Such work is a necessary step towards obtaining robust dynamical BH masses, evaluating the consistency of the main measurement techniques and subsequent effects on the BH scaling relations, and advancing our understanding of the growth of high-mass galaxies and their central BHs.

\TopicNeedSpace{10\baselineskip}
\hypertarget{abs-15}{}
\subsubsection*{Polarimetric signatures of a potential binary supermassive black hole system in Mrk 231.}
\addcontentsline{toc}{subsubsection}{Polarimetric signatures of a potential binary supermassive black hole system in Mrk 231.}
\textbf{Speaker:} Julie Biedermann \\
\textbf{Affiliation:} Université de Strasbourg \textperiodcentered{} Student (BSc/MSc/PhD)\\
Its is believed that there is a supermassive black hole (SMBH) at the center of almost all massive galaxies. When galaxies merge, the SMBHs of the progenitor galaxies are expected to form a bound binary system. Given the high frequency of such mergers throughout cosmic time, binary SMBHs should be relatively common in active galactic nuclei (AGNs). Yet, their direct detection remains a major observational challenge, with most candidates inferred only through indirect signatures.

This presentation will explore how spectropolarimetry can be used to detect and characterize SMBHs binaries in AGNs. We present new perspectives into the internal structure of Markarian 231 (Mrk 231), the nearest ultraluminous infrared galaxy. The optical-UV spectrum of this quasar is distinctive due to a strong attenuation in the near-UV waveband and a localized peak of linear polarization degree associated with a rotation of the polarization angle beyond this ultraviolet break. Our study aims to investigate the hypothesis of a binary SMBHs system in the core of Mrk 231 as proposed by Yan and al (2015) in order to explain the observed enigmatic properties of this object.

To investigate this approach, we reconstruct the total spectral energy distribution (SED) of Mrk 231 by modeling the multi-temperature black body emission of two SMBHs and include the emission of the host galaxy as well as the re-radiation from a dusty torus in the near-infrared. Using this SED model in total flux, we ascribe different values of polarization to each of the components to try to reproduce the unique polarization features of the system across a large wavelength band, ranging from the far-UV to the infrared.

Our total flux SED model based on the binary SMBH hypothesis successfully reproduces the pronounced ultraviolet cut-off observed in the spectrum of Mrk 231, supporting the binary SMBHs scenario. Furthermore, the polarization properties of this quasar can also be reproduced with the same model, using a specific set of parameters that overcomes most of the problems of degeneracy thanks to the simultaneous fitting of the total flux, polarization degree and polarization angle, three independent observables.

This study provides a critical step toward the understanding of the complex structure of Mrk 231 and the role of binary SMBHs in the evolution of AGNs. Such investigations support the case for future high-sensitivity UV spectropolarimeters. In particular, a next-generation instrument like POLLUX, proposed for the Habitable Worlds Observatory, would enable stringent tests of binary SMBH scenarios through features such as double peaks in polarized flux. These observations will provide critical constraints on black hole growth, feedback mechanisms, and the evolution of AGNs.

\TopicNeedSpace{10\baselineskip}
\hypertarget{abs-16}{}
\subsubsection*{WISDOM: A new mm fundamental plane of black hole accretion}
\addcontentsline{toc}{subsubsection}{WISDOM: A new mm fundamental plane of black hole accretion}
\textbf{Speaker:} Martin Bureau \\
\textbf{Affiliation:} University of Oxford \textperiodcentered{} Faculty (tenured)\\
We show how observations of the mm continuum of the central regions of galaxies from the WISDOM project on 10 pc scales reveal a correlation between the nuclear mm continuum luminosities and the supermassive black hole (SMBH) masses, much tighter than most SMBH mass - galaxy correlations. Adding the X-ray luminosities reveals a new "mm fundamental plane of black hole accretion" that is nearly as tight as the SMBH mass - stellar velocity dispersion relation. The relation holds for both low- and high-luminosity active galactic nuclei as well as accreting stellar mass black holes. The relation is reproduced by advection-dominated accretion flow (ADAF) models but not by classic torus-thin accretion disc models. Overall, these new relations thus pose strong constraints on viable accretion disc models, in addition to offering a rapid method to estimate SMBH masses.

\TopicNeedSpace{10\baselineskip}
\hypertarget{abs-17}{}
\subsubsection*{Advancing Black Hole Mass Measurements in Quasars Across the Early Universe}
\addcontentsline{toc}{subsubsection}{Advancing Black Hole Mass Measurements in Quasars Across the Early Universe}
\textbf{Speaker:} Elena Dalla Bonta' \\
\textbf{Affiliation:} University of Padua, Italy \textperiodcentered{} Faculty (tenured)\\
Reliable measurements of black hole masses are essential to reconstruct the assembly history of supermassive black holes and their connection with galaxy evolution. Reverberation mapping remains the most direct technique for active galactic nuclei, yet its demanding observational requirements limit its applicability to large or distant samples. In this contribution, we introduce an advanced single-epoch mass estimation strategy that represents the current state of the art, grounded in the properties of broad emission lines and anchored to a newly homogenized set of reverberation mapping analyses.

The method integrates, in a consistent way, the three physical ingredients that predominantly shape mass determinations: emission line luminosity, velocity width, and accretion rate. Through a systematic reanalysis of archival reverberation mapping campaigns using improved detrending procedures, we establish updated scaling relations that markedly improve the precision of single-epoch estimators. These results can be applied directly to extensive AGN surveys and are particularly timely for present and future spectroscopic programs, including those carried out with JWST.

The typical scatter relative to reverberation based masses is limited to about 0.2–0.3 dex, demonstrating that this approach offers a robust and scalable pathway for investigating black hole populations through cosmic history, opening new perspectives on AGN accretion and galaxy growth from the nearby Universe to the highest redshifts.

\TopicNeedSpace{10\baselineskip}
\hypertarget{abs-18}{}
\subsubsection*{Testing Black Hole-Host Galaxy Scaling Relations in SDSS Close Pairs}
\addcontentsline{toc}{subsubsection}{Testing Black Hole-Host Galaxy Scaling Relations in SDSS Close Pairs}
\textbf{Speaker:} Maggie Huber \\
\textbf{Affiliation:} University of Colorado Boulder \textperiodcentered{} Student (BSc/MSc/PhD)\\
Supermassive black holes (SMBHs) and their host-galaxies co-evolve through channels such as hierarchical merging and AGN feedback, establishing tight scaling relations between SMBH masses and properties of the host galaxy. These scaling relations, particularly those with host galaxy bulge mass (Mbulge) and stellar velocity dispersion (σ), are widely used to estimate the universal SMBH mass function. While there is ongoing investigation into how these different scaling relations predict SMBH mass for isolated galaxies, it is relatively unexplored how scaling relations change for galaxies undergoing a merger. In this talk, I will explore scaling relations in mergers for the first time by taking advantage of the well-established value-added catalogs in the Sloan Digital Sky Survey that provide close galaxy pairs, AGN broad lines, and bulge-disk decomposed stellar masses. We show that differential mass growth between the SMBH and host galaxy during a merger appears to affect the offset between SMBH mass from M•-Mbulge and SMBH mass from AGN broad lines. We find evidence that SMBH mass growth outpaces the bulge for AGN residing in the secondary (less massive) bulge and for smaller physical separations between galaxies in a pair. We also find that M•-Mbulge and M•- σ predict significantly different fractions of major and minor black hole mergers for the full close pairs sample, with M•-Mbulge generally predicting black hole mass ratios closer to 1:1 and more major mergers overall. These results have major implications, especially for predictions of the astrophysical gravitational wave background. In particular, the initial SMBH mass ratio, which we observationally estimate here, is one of the key parameters that is used to determine SMBH binary orbital evolution and detection rates in pulsar timing arrays and LISA.

\TopicNeedSpace{10\baselineskip}
\hypertarget{abs-19}{}
\subsubsection*{Unraveling the broad-line region: novel modeling and deep-learning techniques}
\addcontentsline{toc}{subsubsection}{Unraveling the broad-line region: novel modeling and deep-learning techniques}
\textbf{Speaker:} Kirk Long \\
\textbf{Affiliation:} JILA, University of Colorado Boulder \textperiodcentered{} Student (BSc/MSc/PhD)\\
Active galactic nuclei show considerable line-broadening, an observational hallmark that implies the presence of a supermassive black hole. This broadening implies a characteristic velocity, which, in conjunction with a radial distance measurement—obtained either via reverberation mapping or spectroastrometry—one can estimate the masses of supermassive black holes. But to obtain the most precise masses and unravel the true nature of the broad-line region requires a degree of modeling. We have recently shown that novel models of the BLR are likely required to explain velocity-resolved reverberation data products. In particular, the observed delay and line profiles seem to imply multiple populations of broad line emitting gas with different geometries and kinematics, and we will discuss ongoing efforts to alleviate this tension with multi-component BLR models, enabled by our open-source Julia code BroadLineRegions.jl. We will also showcase a promising deep learning analysis technique for the deconvolution of reverberation mapping data to obtain structural information in the BLR, with applications to similar deconvolution problems in the accretion disk and torus.

Through these approaches we seek to constrain the true nature of the BLR, what physics can plausibly exist there, and how our choice of model systematically influences our measurements of high level properties such as SMBH mass across cosmic time.

\TopicNeedSpace{10\baselineskip}
\hypertarget{abs-20}{}
\subsubsection*{Measuring the Shape of Kerr Black Holes at the Photon Orbit}
\addcontentsline{toc}{subsubsection}{Measuring the Shape of Kerr Black Holes at the Photon Orbit}
\textbf{Speaker:} Kiana Salehi \\
\textbf{Affiliation:} Perimeter Institute \textperiodcentered{} Student (BSc/MSc/PhD)\\
The bright ring-like structures observed in the images of M87* and Sgr A* captured by the Event
Horizon Telescope strongly support the validity of general relativity. Lensed images of the emission
region, often referred to as photon rings in this context, are a direct consequence of the unstable
dynamics of null geodesics near the spherical photon orbit in the Kerr spacetime. The order of the
lensed image can be characterized by the number of half-orbits the photons complete before reaching
the observer, with higher-order photon rings produced by null geodesics that circle the black hole more
times. However, low-order rings are significantly influenced by the astrophysical environment. Mea-
suring the Lyapunov exponent requires probing the exponentially small differences between successive
photon rings or between photon rings and the shadow. We investigate potential astrophysical sources
of systematic error the estimation of Lyapunov exponent, including the location of the observed emis-
sion, and especially at low photon ring order. We show that it is nevertheless possible to measure this
purely gravitational quantity to roughly 10\% and 1\% systematic uncertainty by resolving the n = 2
and n = 3 photon rings with the shadow size, respectively. Therefore, the forthcoming black hole
imaging efforts to capture, even if indirectly, the n= 2 photon ring can result in a measurement of the
Lyapunov exponent that is not limited by astrophysical uncertainties.

\TopicNeedSpace{10\baselineskip}
\hypertarget{abs-21}{}
\subsubsection*{TiDES-RM: Status Updates on the Time-Domain Extragalactic Survey Reverberation Mapping Program}
\addcontentsline{toc}{subsubsection}{TiDES-RM: Status Updates on the Time-Domain Extragalactic Survey Reverberation Mapping Program}
\textbf{Speaker:} Shu Wang \\
\textbf{Affiliation:} University of Southampton \textperiodcentered{} Postdoc\\
The Time-Domain Extragalactic Survey Reverberation Mapping program (TiDES-RM) is a dedicated spectroscopic monitoring campaign with the 4-metre Multi-Object Spectroscopic Telescope (4MOST). TiDES-RM will monitor four LSST deep-drilling fields for at least five years starting in 2026, with a typical cadence of 10–14 days. Combining 4MOST spectroscopy (R\textasciitilde{}4000) with high-cadence LSST photometry, TiDES-RM is designed to measure multi-line reverberation lags and black hole masses for a sample of \textasciitilde{}950 broad-line AGNs over 0.1 < z < 3.5. These data will enable improved single-epoch black hole mass estimators, constraints on BLR kinematics, and tests of the evolution of black hole–host scaling relations across cosmic time. In this talk, I will summarize the program design and sample characterization, present the current observing status, and highlight some of the key science topics that Tides-RM will be working on.

\TopicNeedSpace{10\baselineskip}
\hypertarget{abs-22}{}
\subsubsection*{Weighing Supermassive Black Holes Accurately Across Cosmic Time}
\addcontentsline{toc}{subsubsection}{Weighing Supermassive Black Holes Accurately Across Cosmic Time}
\textbf{Speaker:} Hengyue Zhang \\
\textbf{Affiliation:} University of Oxford \textperiodcentered{} Student (BSc/MSc/PhD)\\
Understanding the mechanisms responsible for the formation and growth of supermassive black holes (SMBHs) and their co-evolution with galaxies requires accurate measurements of SMBH masses across cosmic time. In the local universe, we need a large sample of precise and accurate SMBH mass measurements to constrain the intrinsic scatter in local SMBH-galaxy scaling relations. At high redshifts, we need the first direct dynamical SMBH mass measurements to verify the overmassive nature of early SMBHs. I will first present SMBH mass measurements in local galaxies using ALMA observations of CO kinematics that highly resolve the SMBH’s sphere of influence (SoI), focusing on the measurement in NGC 383. The NGC 383 observation has higher resolution (in physically motivated units) than all other gas-dynamical SMBH mass measurements to date, enabling a 5\%-precision SMBH mass constraint. Additionally, I will present a dynamical constraint on the SMBH mass in the lensed submillimetre galaxy ID141 at z = 4.24, derived by fitting high-resolution ALMA observations directly in visibility space using a novel dynamical and lens forward modelling package.

\TopicNeedSpace{10\baselineskip}
\hypertarget{abs-23}{}
\subsubsection*{Galactic Center Molecular Clouds as Windows into Sgr A*'s Past}
\addcontentsline{toc}{subsubsection}{Galactic Center Molecular Clouds as Windows into Sgr A*'s Past}
\textbf{Speaker:} Stephen DiKerby \\
\textbf{Affiliation:} Michigan State University \textperiodcentered{} Postdoc\\
As the closest supermassive black hole (SMBH) and a target of Event Horizon Telescope, Sgr A* is a unique laboratory for the study of black holes in the local universe. In its quiescent state, Sgr A* emits IR and X-rays via nonthermal processes at about billionth of the Eddington luminosity. However, reflected X-rays from galactic center molecular clouds (GCMC) about 100 light-years from Sgr A* suggest that Sgr A* was in the recent past substantially more luminous than it appears today. New measurements of Iron emission lines by the XRISM telescope tightly constrain the X-ray reflection model while precisely measuring the physical properties of GCMC, enabling detailed modeling of the environment immediately surrounding Sgr A*. Combined with spatially resolved light curves and X-ray polarization measurements, our XRISM observations of GCMC support a new model for Sgr A* as a SMBH that has undergone substantially changes in luminosity on a scale of only a few centuries, revising our understanding of the variability of the only SMBH specimen where these analyses are possible.

\TopicNeedSpace{10\baselineskip}
\hypertarget{abs-24}{}
\subsubsection*{New Determination of Black Hole Masses: Constraints from Reverberation Mapping, Dynamical Modeling, and the Fundamental Plane of the BLR Size–Luminosity–Velocity Relation}
\addcontentsline{toc}{subsubsection}{New Determination of Black Hole Masses: Constraints from Reverberation Mapping, Dynamical Modeling, and the Fundamental Plane of the BLR Size–Luminosity–Velocity Relation}
\textbf{Speaker:} Jong-Hak Woo \\
\textbf{Affiliation:} Seoul National University \textperiodcentered{} Faculty (tenured)\\
Black hole mass is a key property for understanding accretion physics in active galactic nuclei (AGN) and the coevolution of supermassive black holes (SMBHs) with their host galaxies. Reverberation mapping (RM) provides a direct measurement of the size of the broad-line region (BLR), and the black hole mass can be estimated by combining this size with the velocity of the broad emission lines. However, reverberation-based black hole masses carry an uncertainty of \textasciitilde{}0.4 dex, primarily due to the unknown virial factor, which reflects the geometry and inclination of the BLR gas relative to the line of sight.

In this talk, I will first present a comprehensive analysis of BLR kinematics derived from velocity-resolved RM, incorporating results from the Seoul National University AGN Monitoring Project and other major RM campaigns. Using dynamical modeling, we directly determine black hole masses without assuming an average virial factor. The newly derived masses show overall consistency with a modest but non-negligible scatter compared to previous reverberation masses calibrated using an average virial factor based on the M–σ relation.

Second, I will present the fundamental plane relation among BLR size, luminosity, and velocity, based on the best Hβ lag sample of 157 AGNs. We argue that the conventional two-parameter BLR size–luminosity relation is incomplete due to the systematic effect of the Eddington ratio. The BLR fundamental plane provides a new mass estimator, reducing systematic biases in single-epoch black hole mass measurements. We show that masses derived from the previous BLR size–luminosity relation can be overestimated by up to 0.5 dex, implying significant revisions to the cosmic black hole mass density and the inferred growth of black hole seeds in the early universe.


\hypertarget{abs-day-tuesday}{}
\daypalettebanner{Tuesday}{tuesdaybg}{tuesdayaccent}{June 30, 2026}

\section*{Theme 1: Black Hole Demographics and Scaling Relations}
\addcontentsline{toc}{section}{Theme 1: Black Hole Demographics and Scaling Relations}

\TopicNeedSpace{10\baselineskip}
\hypertarget{abs-25}{}
\subsubsection*{Strong lensing of AGN as a cosmic probe of SMBH structure and multiplicity at high redshift.}
\addcontentsline{toc}{subsubsection}{Strong lensing of AGN as a cosmic probe of SMBH structure and multiplicity at high redshift.}
\textbf{Speaker:} Schwartz Dan \\
\textbf{Affiliation:} Harvard University\\
Strong gravitational lensing of distant AGN offers a powerful cosmic telescope to probe sub‑kiloparsec structures at early cosmic epochs when spatial offsets and mergers are expected to be abundant, yet unresolvable by direct imaging with current and future instruments. I will present an innovative solution to this conundrum, combining lens mass modeling in the optical with spatially resolved Chandra X‑ray imaging of lensed AGN. Chandra’s sub-arcsecond angular resolution is critical for resolving individual lensed images in the X-rays. By performing a maximum-likelihood analysis, we recover the relative spatial origin of the X‑ray emission with respect to the optical, achieving milliarcsecond relative astrometry when using Gaia observations. 
I will showcase recent results from Chandra observations of three lensed quasars: J0659 (z = 3.1), J1433 (z = 2.74), and the first zig-zag lens J1721 (z = 2.37), highlighting sensitivity to structure at tens-to-hundreds of parsecs expected from low-z AGN studies. I will then discuss prospects for identifying dual AGN at sub‑kpc separations at high-z – critical predictions in galaxy formation, yet observationally elusive. I will introduce our current multiwavelength results addressing the nature of the multiply-imaged radio-loud system J2019 (z = 3.3), a strong dual AGN candidate at a projected separation of 175 pc. If confirmed, this would challenge galaxy merger rates assumed by state-of-the-art cosmological models.
Finally, I will describe how the caustic method will scale to the thousands of new lensed AGN expected from Euclid, LSST, and other discovery engines, leveraging machine-learning-optimized lens modeling for sample-level analyses.

\section*{Theme 4: First Black Holes, Seeds and High-redshift}
\addcontentsline{toc}{section}{Theme 4: First Black Holes, Seeds and High-redshift}

\TopicNeedSpace{10\baselineskip}
\hypertarget{abs-5}{}
\subsubsection*{Observational constraints on distant black holes and their immediate environments}
\addcontentsline{toc}{subsubsection}{Observational constraints on distant black holes and their immediate environments}
\textbf{Speaker:} Hannah Ubler \invitedbadge{}\\

One of the most surprising JWST results is an unexpectedly high abundance of actively accreting black holes in the early Universe. Compared to the local population, these black holes differ in various aspects, such as their relation to their host galaxies or their multi-wavelength properties. These observational findings challenge our understanding of the past evolution of present-day supermassive black holes, and provide new ways to constrain theoretical models of black hole formation and growth. I will give an overview of recent observational results on massive black holes in the first few billion years, driven by the unprecedented capabilities of JWST to explore cosmic dawn, and especially by NIRSpec-IFS to constrain early black holes and their immediate environments.

\TopicNeedSpace{10\baselineskip}
\hypertarget{abs-6}{}
\subsubsection*{High-redshift black hole formation and growth: a theoretical perspective}
\addcontentsline{toc}{subsubsection}{High-redshift black hole formation and growth: a theoretical perspective}
\textbf{Speaker:} Sophie Koudmani \invitedbadge{}\\

Supermassive black holes (SMBHs) reside at the centres of most massive, if not all, galaxies influencing their host galaxy’s evolution through a complex feedback cycle. SMBH feeding and feedback spans a vast range of scales, from the accretion disc to the cosmic web, so that all past cosmological simulations have had to rely on simple ‘subgrid’ models for SMBH evolution, limiting their predictive capabilities. Recent observational breakthroughs have challenged these models as JWST has detected significantly more active SMBHs in the early Universe than had been predicted. The recent evidence for a gravitational wave background from SMBH mergers also points to the need for more efficient early SMBH growth. The origin of SMBHs is a major open research question due to long-standing degeneracies between SMBH seed masses and growth rates. I will discuss the latest simulation suites that aim to uncover the formation and rapid growth of SMBHs in the early Universe, particularly in light of new observations from JWST. I will outline the main challenges, including uncertainties for the seeding channels and the AGN accretion modelling as well as, crucially, mimicking observational selection functions for apple-to-apple comparisons between the simulations and observations. Finally, I will conclude with an outlook on next-generation SMBH models and synthetic observation tools for galaxy formation simulations.

\TopicNeedSpace{10\baselineskip}
\hypertarget{abs-10}{}
\subsubsection*{Little Red Dots: Dense Gas, Broad Lines, and the Nature of a New AGN Population}
\addcontentsline{toc}{subsubsection}{Little Red Dots: Dense Gas, Broad Lines, and the Nature of a New AGN Population}
\textbf{Speaker:} Francesco D'Eugenio \invitedbadge{}\\

Little Red Dots (LRDs) are a newly identified class of compact, broad-line AGN at high redshifts whose high number density suggests they play an important role in early black-hole growth. They are characterized by broad permitted lines, V-shaped SEDs, and a remarkable incidence of unusual spectral features compared to quasars and local AGN. They show Balmer-line absorption at a rate 10-100 times higher than in normal AGN, implying the presence of gas with high densities and large column densities. Many also exhibit metal-line absorption, strong Balmer breaks, low-ionization emission, exponential line wings, and, in some cases, forbidden iron and strong auroral emission, alongside a general paucity of high-ionization lines.

I will discuss a spectroscopic analysis of the absorbing and emitting gas in LRDs, highlighting its complex kinematic structure and its likely connection to both feeding and feedback around low-mass black holes. I will discuss these phenomena in the context of density stratification, gas clouds or stationary envelopes surrounding the central engine, and the broader physical picture emerging for this enigmatic population.

\TopicNeedSpace{10\baselineskip}
\hypertarget{abs-54}{}
\subsubsection*{Evolution of supermassive black holes and their host galaxies at z > 6 with JWST}
\addcontentsline{toc}{subsubsection}{Evolution of supermassive black holes and their host galaxies at z > 6 with JWST}
\textbf{Speaker:} John Silverman \\
\textbf{Affiliation:} Kavli IPMU; University of Tokyo \textperiodcentered{} Faculty (tenured)\\
Clues on the formation of supermassive black holes may be found in the properties of their host galaxies and mass relations with cosmic time. We will present results from JWST programs, while focusing on the Subaru High-z Exploration of Low-Luminosity Quasars (SHELLQs) survey at z > 6. After image and spectral decomposition of NIRCam and NIRSpec data, host detections reveal their stellar masses, morphology, and properties of the stellar population. Coupled with black hole masses, we carry out one of the first assessments of the intrinsic ratio between black hole mass and stellar mass at z > 6. Lastly, the connection with lower mass black holes, including the “Little Red Dot” population, presents new insights into early phases of black hole growth and their seeding.

\TopicNeedSpace{10\baselineskip}
\hypertarget{abs-55}{}
\subsubsection*{A Multi-Wavelength View of z>6 QSOs: Combining the deepest X-ray Observations with JWST/NIRSpec Spectroscopy}
\addcontentsline{toc}{subsubsection}{A Multi-Wavelength View of z>6 QSOs: Combining the deepest X-ray Observations with JWST/NIRSpec Spectroscopy}
\textbf{Speaker:} Luca Zappacosta \\
\textbf{Affiliation:} INAF - OAR \textperiodcentered{} Research staff / Scientist\\
Shedding light on the nuclear properties of z>6 quasars (QSOs) powered by supermassive black holes (SMBHs) with masses >10\textasciicircum{}9 Msun is fundamental to understand their rapid formation. We probe the innermost regions of these QSOs via X-ray emission which, at just a few tens of gravitational radii, provides the most direct window into the growing SMBH.

We present results from HYPERION, an unprecedented \textasciitilde{}700-hour XMM-Newton Heritage Program targeting 18 luminous QSOs at z=6-7.5 selected for their extreme SMBH growth. Our findings reveal a strikingly new regime for the nuclear properties of these early-formed QSOs: their X-ray continuum slopes (Γ) are significantly steeper than those of analogous QSOs at z<6. By investigating correlations between Γ and UV-derived properties, we identify a significant link between the X-ray slope and the velocity of nuclear AGN-driven ionized winds, potentially indicative of fast SMBH accretion.

This investigation is further expanded by a new 150-hour XMM-Newton Large Program which, combined with archival X-ray data, brings the total sample to 30 QSOs and includes sources with slower growth rates. The data confirm that steep X-ray slopes are a ubiquitous feature of the z>6 QSOs population, independent of growth-rate selection. We discuss the implications of these results to explain the widespread non-detection of X-rays in the low-luminosity AGN population recently discovered by JWST.

Finally, we report on the ongoing JWST/NIRSpec spectral analysis of more than half of this sample. By characterizing the rest-frame optical emission features, we provide the first systematic investigation of the broad/narrow-line region and large-scale ionized outflows, and their link with the innermost X-ray emitting corona, offering an unprecedented, multi-component view of the most distant QSOs currently known.

\TopicNeedSpace{10\baselineskip}
\hypertarget{abs-56}{}
\subsubsection*{A Population of Local Red Dots in SDSS}
\addcontentsline{toc}{subsubsection}{A Population of Local Red Dots in SDSS}
\textbf{Speaker:} Quinn Casey \\
\textbf{Affiliation:} Dartmouth College \textperiodcentered{} Student (BSc/MSc/PhD)\\
Compact and red sources in the high redshift (z\textasciitilde{}5) Universe, known as "Little Red Dots" (LRDs), are among JWST's most intriguing discoveries. These sources have broad Balmer emission lines, weak X-ray emission, and unique spectral energy distributions (SEDs) poorly fit by either stellar or AGN templates. Local analogs of LRDs allow for detailed studies of the underlying physical processes with archival multi-wavelength datasets unavailable in the high-z Universe. We show that the SDSS ugriz filters at z=0.4, 0.8 overlap well with the JWST filters used to select LRDs at z\textasciitilde{}5. We define a sample of \textasciitilde{}1300 Local Red Dots (LoRDs) from SDSS quasars which share the same photometric colors of LRDs. A subset of the LoRD sample selected to have a V-shaped continua show prominent higher-order Balmer absorption features and [NeV]3426 emission, both of which would be missed in JWST/PRISM observations given the low spectral resolution. A composite SED of the LoRDs differs from a typical quasar SED in the UV/optical, but the two agree with each other in the NIR. The LoRD SED matches well with a stack of LRDs and can be modeled with a reddened AGN, a hot blackbody, and a variety of host galaxy templates. Interestingly, the LoRDs are detected in the X-rays at a rate comparable to the control sample. Whether this is due to different properties between the LoRDs and the LRDs, or the LoRDs having a soft X-ray spectrum unobservable at high-z is a question we aim to answer.

\TopicNeedSpace{10\baselineskip}
\hypertarget{abs-57}{}
\subsubsection*{On the Evolution of Massive Black Holes and AGN in High-Redshift, Bursty Galaxies and Insights into the Nature of Little Red Dots}
\addcontentsline{toc}{subsubsection}{On the Evolution of Massive Black Holes and AGN in High-Redshift, Bursty Galaxies and Insights into the Nature of Little Red Dots}
\textbf{Speaker:} Claude-André Faucher-Giguère \\
\textbf{Affiliation:} Northwestern University \textperiodcentered{} Faculty (tenured)\\
High-resolution galaxy formation simulations predict that high-redshift, low-mass galaxies undergo bursty star formation, with star formation rates that vary at the order-of-magnitude level on short timescales. Star formation bursts repeatedly eject gas from galaxies, strongly limiting the gas available to accrete onto black holes. Moreover, the clumpy, highly dynamic morphologies of high-redshift galaxies imply they often do not have well-defined dynamical centers toward which black holes can efficiently sink, presenting an additional challenge — beyond the short age of the universe — in explaining the existence of massive nuclear black holes during the first gigayear of cosmic history. In this talk, I will summarize results from several theoretical studies based on the FIRE galaxy formation simulations that provide insights into the high-redshift AGN population and the origin of supermassive black holes. These results are grounded in simulations that are in excellent agreement with many observed properties of high-redshift host galaxy populations, including their UV luminosity function and metallicities measured by the James Webb Space Telescope. As time permits, I will discuss the orbits of seed black holes in early galaxies and the possible role of star clusters, predictions for the AGN luminosity function and how it is affected by bursty star formation, and how the recently discovered population of Little Red Dots can be explained in the context of nuclear accretion in realistic bursty galaxies.

\TopicNeedSpace{10\baselineskip}
\hypertarget{abs-58}{}
\subsubsection*{Peeling Back Little Red Dots with JWST/NIRSpec IFU: Spatially Resolving a Blue Host and Compact Black Hole Star}
\addcontentsline{toc}{subsubsection}{Peeling Back Little Red Dots with JWST/NIRSpec IFU: Spatially Resolving a Blue Host and Compact Black Hole Star}
\textbf{Speaker:} Yuzo Ishikawa \\
\textbf{Affiliation:} Massachusetts Institute of Technology \textperiodcentered{} Postdoc\\
Little Red Dots (LRDs) are a population of compact, faint, and extremely red sources recently identified by JWST. To date, broadband imaging and slit/grism spectroscopy have shown that LRDs exhibit a mysteriously complex spectrum: a “V-shaped” blue/red continuum, narrow emission lines, broad Balmer emission lines, and in some cases Balmer absorption. While the physical origin of these components remains debated, emerging studies propose that they arise from a compact accreting black hole, surrounded by dense gas, within a host galaxy. We test this model with deep, JWST/NIRSpec IFU observations of 5 broad Hα-selected LRDs using the PRISM and G395H modes. By jointly decomposing the spatial and spectral properties of the continuum, emission lines, and absorption components, we directly separate LRDs into two physically distinct structures. We detect extended blue continuum that is co-spatial with the extended narrow-line emission (e.g., [O III] and Hα), establishing the presence of a blue host galaxy. In contrast, the red continuum is compact and nearly co-spatial with the broad Balmer emission and absorption, consistent with the “Black Hole Star” model. These results provide direct, spatially resolved evidence that the narrow- and broad-line components arise from physically distinct regions. Together, our findings clarify the internal structure of LRDs and place new observational constraints on their role in early galaxy and massive black hole evolution.

\TopicNeedSpace{10\baselineskip}
\hypertarget{abs-60}{}
\subsubsection*{Setting Little Red Dots in FIRE: The Ability of Bursty Galaxies to Host an Abundant Population of High-z AGN}
\addcontentsline{toc}{subsubsection}{Setting Little Red Dots in FIRE: The Ability of Bursty Galaxies to Host an Abundant Population of High-z AGN}
\textbf{Speaker:} Andy Marszewski \\
\textbf{Affiliation:} Northwestern University / CIERA \textperiodcentered{} Student (BSc/MSc/PhD)\\
The James Webb Space telescope has unveiled an abundant population of potential active galactic nuclei (AGN) at high redshift known as little red dots (LRDs), which are likely hosted in relatively low-mass galaxies. However, previous theoretical models have highlighted the difficulty in continuously feeding massive black holes in the central regions of bursty, high-redshift, low-mass galaxies because of repeated gas evacuation by stellar feedback. We analyze galaxies within high-redshift FIRE-2 simulations to understand whether they are capable of hosting the observed abundant population of high-redshift AGN. We use the gravitational torque-driven accretion model and a simple free-fall accretion model to derive black hole accretion rates and construct predicted AGN bolometric luminosity functions for z=5-7. Despite the highly intermittent black hole accretion rates inferred from our FIRE-2 sample, we show that bursty, high-redshift galaxies are predicted to host a sufficiently abundant population of AGN in the luminosity range of LRDs. Interestingly, our models reproduce the observed number of AGN at the highest LRD luminosities but overpredict their abundance at lower luminosities relative to observations. We explore potential explanations for this overprediction and their implications for the types of galaxies that host LRDs and for black hole growth in the early universe. These models also predict that the vast majority of AGN in the luminosity range of observed LRDs are hosted by low-mass (log⁡(M\_⋆/M\_⊙ )=5-10) galaxies.

\TopicNeedSpace{10\baselineskip}
\hypertarget{abs-61}{}
\subsubsection*{The rapid formation and growth of supermassive black hole seeds in dense star clusters at high redshift}
\addcontentsline{toc}{subsubsection}{The rapid formation and growth of supermassive black hole seeds in dense star clusters at high redshift}
\textbf{Speaker:} Thorsten Naab \\
\textbf{Affiliation:} Max Planck Institute for Astrophysics \textperiodcentered{} Faculty (tenured)\\
JWST reveals massive, compact star clusters at high redshift as well as early supermassive black holes. We test whether compact star clusters can produce the seeds of supermassive black holes using high-accuracy direct N-body simulations with the novel BIFROST code. The simulations follow the hierarchical assembly of clusters with up to several million stars, including stellar evolution, initial stellar multiples, and post-Newtonian dynamics for black-hole encounters. We find that intermediate-mass black holes (IMBHs) of up to several thousand solar masses form rapidly, within a few million years, through runaway collisions of massive stars. Subsequent tidal disruption events and gravitational wave driven mergers among IMBHs - predicted by this channel - further grow the remnants. A large parameter study shows that this pathway is most efficient in dense, low-metallicity clusters typical of the early Universe. We discuss implications for the cosmic IMBH formation rate and show that gas accretion can accelerate growth when the IMBHs are embedded in a nuclear star cluster, providing viable seeds for early supermassive black holes.

\TopicNeedSpace{10\baselineskip}
\hypertarget{abs-62}{}
\subsubsection*{Twinkle, twinkle, little pair: the tightest z\textasciitilde{}6 quasar duo as a beacon of early structure formation}
\addcontentsline{toc}{subsubsection}{Twinkle, twinkle, little pair: the tightest z\textasciitilde{}6 quasar duo as a beacon of early structure formation}
\textbf{Speaker:} Silvia Onorato \\
\textbf{Affiliation:} Gemini Observatory \textperiodcentered{} Postdoc\\
The discovery of close quasar pairs at z \textasciitilde{} 6 provides a new avenue to investigate early supermassive black hole (SMBH) growth, galaxy assembly, and the small-scale clustering of quasars near the end of cosmic reionization. We present the discovery and first characterization of the most compact known quasar pair in the Epoch of Reionization, J1602+4228, with a projected separation of only \textasciitilde{}0.8’’ (\textasciitilde{} 4.5 kpc at z \textasciitilde{} 6.08). This system offers a rare opportunity to probe SMBH fueling and early structure formation on kiloparsec scales at high redshift. Using spectroscopy from Gemini/GNIRS and JWST/NIRSpec IFU, HST/WFC3 imaging, and [C II] maps from IRAM NOEMA, we characterize the emission-line properties, host-galaxy morphology, and dynamical context of the two quasars, enabling constraints on their black hole masses, accretion states, and potential signs of interaction or merger-driven activity. By comparing the environment of J1602+4228 to expectations from models of early quasar clustering, we assess whether such close pairs preferentially inhabit overdense regions or can arise in more typical halos. This work represents a first step toward establishing the physical drivers of SMBH growth during the EoR and the role of early mergers and environment in shaping the first massive black holes.

\TopicNeedSpace{10\baselineskip}
\hypertarget{abs-63}{}
\subsubsection*{Little Red Dots: one observational tag, diverse spectroscopic flavors with distinct stellar and nuclear activity features}
\addcontentsline{toc}{subsubsection}{Little Red Dots: one observational tag, diverse spectroscopic flavors with distinct stellar and nuclear activity features}
\textbf{Speaker:} Pablo G. Pérez-González \\
\textbf{Affiliation:} Centro de Astrobiología (CAB), CSIC-INTA \textperiodcentered{} Faculty (tenured)\\
We will present an spectro-photometric analysis of a complete sample of little red dots (LRDs). We stack NIRSpec prism data for \textasciitilde{}250 source, and add MIRI observations to produce high SNR (>30) spectral energy distributions of LRDs covering the rest-frame wavelength range from 1000 Angstrom to 4 micron. Our analysis concentrates in some spectral features of stellar emission present in all LRDs, more interestingly, Wolf Rayet signatures. In addition, we analyze iron emission that we use to estimate accretion properties, finding relatively large values but sub-Eddington. In addition, we present a battery of line ratios that characterize the gas content and structure of LRDs, challenging quasi-star scenarios. Based on a comprehensive analysis of the spectrophotometric data, we discuss the validity of black hole and stellar masses, finding that LRDs host moderately over-massive black holes, in contrast with other estimations based on single-epoch line emission.

\TopicNeedSpace{10\baselineskip}
\hypertarget{abs-64}{}
\subsubsection*{Observing the First Supermassive Black Holes with JWST}
\addcontentsline{toc}{subsubsection}{Observing the First Supermassive Black Holes with JWST}
\textbf{Speaker:} Hyun Song \\
\textbf{Affiliation:} Seoul National University \textperiodcentered{} Student (BSc/MSc/PhD)\\
Authors: Hyun Song, Junehyoung Jeon, Anthony Taylor, Volker Bromm

The James Webb Space Telescope (JWST) has significantly expanded the frontier of observational cosmology, and the recent discovery of supermassive black holes (SMBHs) at z \textasciitilde{}10 promises to directly probe their formation and early evolution. Such observations may help to answer how SMBHs were born and grew into the billion-solar-mass objects observed at the centers of quasars and local galaxies. Here, we investigate the redshift horizon at which JWST can detect the first SMBHs. This horizon is determined by balancing the minimum black hole mass required to produce sufficient flux above the JWST sensitivity threshold and the maximum mass allowed by the abundance of SMBHs. We examine three SMBH formation and growth scenarios: a feedback-regulated model, a light-seed model accreting at twice the Eddington limit, and a direct-collapse black hole (DCBH) seed accreting at the Eddington limit.
 We find that the feedback-regulated model predicts a redshift horizon of z\textasciitilde{}7, which is already exceeded by current JWST observations, indicating that feedback-regulated growth alone cannot explain the earliest SMBHs. In contrast, the light seed and DCBH scenarios yield redshift horizons of z\textasciitilde{}15 and z\textasciitilde{}17, respectively. We further show that spectroscopic confirmation of an SMBH at these redshifts may be feasible through detection of strong UV emission lines. This suggests promising prospects for pushing observational studies of primitive SMBHs to epochs as early as \textasciitilde{}250 million years after the Big Bang. Moreover, our results imply that detection of an SMBH at z > 15 would strongly favor a heavy seed origin.

\TopicNeedSpace{10\baselineskip}
\hypertarget{abs-66}{}
\subsubsection*{The Implication of High-z Accretion for SMBH Jet Physics and Seed Environments}
\addcontentsline{toc}{subsubsection}{The Implication of High-z Accretion for SMBH Jet Physics and Seed Environments}
\textbf{Speaker:} Greg Bryan \\
\textbf{Affiliation:} Columbia University \textperiodcentered{} Faculty (tenured)\\
Recent observations suggest that some high-redshift black holes grow at or above the Eddington limit, posing a challenge for models in which mechanical feedback strongly self-regulates accretion. I will discuss a framework that connects small-scale jet physics to cosmological black hole assembly in order to assess what such rapid growth implies for feedback efficiencies and host environments. High-resolution simulations show that jet-driven outflows can efficiently disrupt gas in low-mass halos, limiting growth unless the mechanical coupling is weak or the ambient gas density is unusually high. When embedded in cosmological models, these results indicate that early black hole evolution is highly sensitive to feedback prescriptions and halo structure. If sustained super-Eddington growth is common, it may point to inefficient mechanical feedback and/or seed formation in halos more massive and centrally concentrated than the standard atomic-cooling picture.


\hypertarget{abs-day-wednesday}{}
\daypalettebanner{Wednesday}{wednesdaybg}{wednesdayaccent}{July 1, 2026}

\section*{Theme 3: Accretion, Jets and Multi-messenger Signatures}
\addcontentsline{toc}{section}{Theme 3: Accretion, Jets and Multi-messenger Signatures}

\TopicNeedSpace{10\baselineskip}
\hypertarget{abs-4}{}
\subsubsection*{Feeding The Dragons: Observational Perspective}
\addcontentsline{toc}{subsubsection}{Feeding The Dragons: Observational Perspective}
\textbf{Speaker:} Megan Donahue \invitedbadge{}\\

We don’t know exactly how black holes adjust their fueling rates to the large scale halo properties of their galaxy and cluster hosts. But
observations show that’s what they have done by about redshift 1. Observations also yield clues as to how they do it. I will endeavor
to review what we know, what we think we know, and what we’d like to find out based multiwavelength, multiscale observations of supermassive black holes and their baryon halo.

\TopicNeedSpace{10\baselineskip}
\hypertarget{abs-33}{}
\subsubsection*{Multiphase Radio-Mode AGN Feedback as Seen by JWST}
\addcontentsline{toc}{subsubsection}{Multiphase Radio-Mode AGN Feedback as Seen by JWST}
\textbf{Speaker:} Julie Hlavacek-Larrondo \\
\textbf{Affiliation:} Université de Montréal \textperiodcentered{} Faculty (tenured)\\
The best place to study radio-mode AGN feedback is in the hot atmospheres of galaxy clusters, which host the most massive black holes and where we can directly image their impact on the surrounding medium. At cluster centers, supermassive black holes launch powerful relativistic jets that inject vast amounts of energy through shock fronts, sound waves, turbulence, and molecular outflows. In this talk, I will present novel multi-wavelength observations that, for the first time, allow us to trace the entire radio-mode feedback cycle—from heating, to cooling, to feeding, and across scales from kiloparsecs down to parsecs. First, I will present new JWST observations of the archetypal Centaurus cluster, where NIRSpec/IFU data reveal a rotating circumnuclear disk directly connected to the surrounding filament network. This structure provides the long-sought missing link that channels gas from kiloparsec to 100-pc scales and ultimately into the central black hole, effectively closing the feedback loop. I will then place these results in a broader context using the Perseus cluster, where the combination of JWST with state-of-the-art XRISM and SITELLE/IFU observations allows us to map the full multiphase feedback cycle. In this system, I will show that the hot X-ray–emitting gas is deeply intertwined with the cold and warm nebular components, demonstrating that AGN feedback is inherently multiphase and dynamically complex. Together, these findings shed new light on the fundamental mechanisms that drive galaxy evolution, particularly supermassive black hole fueling/feedback and its intricate interplay with the host galaxy.

\TopicNeedSpace{10\baselineskip}
\hypertarget{abs-34}{}
\subsubsection*{Inside the Belly of the Beast: A Sharp New Look at AGN Feedback with AO-Assisted MUSE}
\addcontentsline{toc}{subsubsection}{Inside the Belly of the Beast: A Sharp New Look at AGN Feedback with AO-Assisted MUSE}
\textbf{Speaker:} Stéphanie Juneau \\
\textbf{Affiliation:} NSF NOIRLab \textperiodcentered{} Research staff / Scientist\\
NGC 7582 is a nearby galaxy with a heavily obscured active galactic nucleus (AGN). Prior MUSE integral-field spectroscopy revealed a 600 parsec circumnuclear star-forming ring and a striking biconical ionized outflow extending over kiloparsec scales. This inner ring likely helps collimate the AGN-driven wind and contributes to the heavy nuclear obscuration. Using MUSE with adaptive optics, we resolve the central region at <0.1arcsec (10 parsec) spatial resolution. This unprecedented view reveals a previously unseen compact ionized outflow (10 pc across) embedded within the known large-scale outflow bicone. The small-scale outflow originates in the immediate vicinity of the AGN, carrying high-velocity gas that was blurred in earlier data. It lies fully inside the footprint of the kpc-scale outflow, indicating a nested, multi-scale wind structure. The discovery of this inner outflow demonstrates that AGN feedback in NGC 7582 operates on nested scales from 10 pc to over 1 kpc. Our findings underscore the importance of high-resolution IFU observations in linking the nuclear wind launching region to its galaxy-scale impact. NGC 7582 now serves as a striking example of how an obscured, highly accreting AGN can drive multi-scale outflows and interact with host galaxy substructure as part of the feedback process. While such studies are currently limited to nearby or lensed galaxies, they preview the capabilities of future 3D spectroscopic observations at high redshift with JWST and ELTs.

\TopicNeedSpace{10\baselineskip}
\hypertarget{abs-38}{}
\subsubsection*{Central star formation in local AGN hosts: is there tension between coeval growth and AGN feedback?}
\addcontentsline{toc}{subsubsection}{Central star formation in local AGN hosts: is there tension between coeval growth and AGN feedback?}
\textbf{Speaker:} Qingling Ni \\
\textbf{Affiliation:} MPE \textperiodcentered{} Postdoc\\
While negative AGN feedback may suppress star formation in galaxy centers, the coeval growth scenario of black holes and galaxies predicts a close observational link between AGN activity and central star formation. Utilizing a sample of approximately 3000 galaxies in the Mapping Nearby Galaxies at Apache Point Observatory (MaNGA) survey that have X-ray coverage,  we investigate how AGN fraction depends on the star formation rate surface density within the central 1 kpc region (ΣSFR,1 kpc). We find that the fraction of X-ray AGN with relatively high specific black hole accretion rates increases with ΣSFR,1 kpc, consistent with the coeval growth picture. Comparison of the mean star formation rate surface density (ΣSFR) profiles of the host galaxies of these X-ray AGN and matched normal galaxies reveals elevated ΣSFR in X-ray AGN hosts across the entire central region. As for optically-selected AGN, their hosts also tend to show elevated ΣSFR in the central regions, but differ from X-ray AGN in terms of the trend of AGN fraction vs. ΣSFR,1 kpc, which can be explained by selection effects. Overall, these trends support the coeval growth picture, and they also do not contradict observational evidence for AGN feedback, as the time-averaged impact of local AGN-driven outflows in our sample appears to be modest.

\TopicNeedSpace{10\baselineskip}
\hypertarget{abs-40}{}
\subsubsection*{Ruling Out Compact Jets as the Dominant Source of Radio Emission in Radio-quiet, High-Eddington-ratio Active Galactic Nuclei}
\addcontentsline{toc}{subsubsection}{Ruling Out Compact Jets as the Dominant Source of Radio Emission in Radio-quiet, High-Eddington-ratio Active Galactic Nuclei}
\textbf{Speaker:} Jeremiah Paul \\
\textbf{Affiliation:} Texas Tech University \textperiodcentered{} Postdoc\\
The question of whether active galactic nuclei (AGNs) produce jets at high accretion rates holds powerful implications for our understanding of how early seed black holes grew. Intriguingly, general relativistic magnetohydrodynamics simulations often predict the launching of moderately large-scale jets from super-Eddington accretion flows, but this prediction seems at odds with observations indicating most high/super-Eddington AGNs are radio quiet. Core radio emission in radio-quiet AGNs can originate from multiple processes with varied spectral signatures and morphologies, but the dominant emission mechanism is still actively debated. Apart from compact jets, likely sources include coronal activity, outflows/winds and shocks, and star formation activity. We used the ratio of radio to X-ray luminosities to probe the dominant source of radio emission for a sample of 69 radio-quiet, high/super-Eddington AGNs with black-hole masses \textasciitilde{} 10\textasciicircum{}5 - 10\textasciicircum{}9 M\_sun. With this wide dynamic range in black-hole masses, we adapted existing formalisms for how jet-dominated radio emission and accretion-powered X-ray emission scale with black-hole mass into the super-Eddington regime (where an AGN is expected to enter a radiatively inefficient 'slim disk' state). Here, I show how the radio/X-ray luminosity ratios observed across this mass range appear inconsistent with a jet-dominated model for radio emission. I discuss how our results may instead suggest a combination of coronal and outflow activity, with increased excess from strong outflows at higher accretion rates. I also discuss limitations of the present study and explore future opportunities to observe the mechanisms for radio emission and better understand how much accretion power may be advected directly into the black hole vs. radiated away or ejected by outflows/jets.

\TopicNeedSpace{10\baselineskip}
\hypertarget{abs-41}{}
\subsubsection*{Probing the Physical State of Warm Gas in NGC 4696: New Infrared Diagnostics for JWST/NIRSpec}
\addcontentsline{toc}{subsubsection}{Probing the Physical State of Warm Gas in NGC 4696: New Infrared Diagnostics for JWST/NIRSpec}
\textbf{Speaker:} Olivia Pereira \\
\textbf{Affiliation:} Université de Montréal \textperiodcentered{} Student (BSc/MSc/PhD)\\
The best place to study AGN feedback processes is in galaxy clusters, with Brightest Cluster Galaxies (BCGs) serving as ideal test cases for observations. While the macro-scale effects of AGN feedback are well documented, the physical mechanisms that excite the warm molecular gas at sub-kiloparsec scales, required to bridge the gap between the hot ICM and the supermassive black hole (SMBH) accretion flow, remain a subject of intense debate. 
In this talk, I will present for the first time Cycle 3 JWST/NIRSPec Integral Field Unit (IFU) observations of the innermost 640 x 640 pc region of NGC 4696, the prototypical cool-core BCG of the Centaurus cluster. I will report the robust detection of over 20 atomic and molecular lines, with emphasis on the rovibrational H₂ S-series transitions. These lines demonstrate significant spatial variation across the map, revealing structures previously blurred in ground-based data. Based on Cloudy simulations, I will present novel infrared versions of the BPT diagram optimized for the JWST era. Using this framework, I will provide the first interpretations of the gas’s physical state across distinct sub-regions, resolved at previously inaccessible scales. This work highlights the transformative power of JWST in quantifying the energetics of feeding and feedback processes at the heart of the most massive galaxies in the local universe.

\TopicNeedSpace{10\baselineskip}
\hypertarget{abs-42}{}
\subsubsection*{News from the Bright End: Extreme Outflows and X-ray Weakness in Luminous Quasars from WISSH}
\addcontentsline{toc}{subsubsection}{News from the Bright End: Extreme Outflows and X-ray Weakness in Luminous Quasars from WISSH}
\textbf{Speaker:} Enrico Piconcelli \\
\textbf{Affiliation:} Osservatorio Astronomico di Roma (INAF) \textperiodcentered{} Research staff / Scientist\\
Highly luminous quasars provide the best laboratories to investigate the bright end of feedback and accretion processes. In this talk, I will present recent multi-wavelength results from the WISSH survey, which targets a sample of \textasciitilde{}100 quasars with bolometric luminosities Lbol> 10\textasciicircum{}47 erg/s at Cosmic Noon.
We find that a significant fraction of WISSH QSOs exhibit ultra-fast (v > 0.1c) BAL outflows in the UV, with high kinetic power sufficient to drive efficient feedback. Multi-epoch spectroscopy of a subset of these extreme BAL quasars reveals decelerating outflows, likely caused by interactions with the interstellar medium, which may provide a rare demonstration of feedback in action.
Despite the narrow range in Lbol, WISSH QSOs show a broad dispersion in X-ray luminosity, with about one-third classified as intrinsically X-ray weak. This indicates significant heterogeneity in the X-ray corona and inner accretion flow, both within this population and compared to typical, less powerful AGN. Finally, the blueshift of the CIV emission line, tracing nuclear ionized winds, strongly depends on X-ray luminosity but not on UV or NIR luminosities. These results suggest a close coupling between the X-ray corona and the kinematics of the inner emission-line region in these highly accreting, powerful AGN.

\TopicNeedSpace{10\baselineskip}
\hypertarget{abs-49}{}
\subsubsection*{A JWST/MIRI-MRS view of the Galactic Center : new evidence for recent star formation}
\addcontentsline{toc}{subsubsection}{A JWST/MIRI-MRS view of the Galactic Center : new evidence for recent star formation}
\textbf{Speaker:} Pierre Vermot \\
\textbf{Affiliation:} LIRA - Observatoire de Paris \textperiodcentered{} Postdoc\\
We present new JWST/MIRI-MRS observations of the central parsec of the Galactic Center, targeting the circumnuclear disk (CND) and the central cavity (CC) around the supermassive black hole. Using full 5–27 μm spectroscopy combined with CLOUDY photoionization modeling, we decompose the mid-infrared emission into distinct gas phases and map their spatial distribution.

In both regions, the emission is dominated by warm ionized gas at temperatures of 10⁴–10⁴.⁸ K. Molecular gas is prominent and structured in the CND but largely absent in the CC, while coronal gas is detected at the interfaces between phases. An elongated, hot coronal feature is identified in the CND, oriented perpendicular to the direction toward the SMBH, suggestive of a large-scale shock.

Abundance measurements reveal enhanced α and CNO elements with strong Fe depletion, consistent with recent enrichment by core-collapse supernovae and massive stellar winds, and a minimal contribution from Type Ia supernovae. This chemical pattern provides new evidence for recent massive star formation in the immediate vicinity of the Galactic Center, rather than long-term cumulative enrichment.

\TopicNeedSpace{10\baselineskip}
\hypertarget{abs-52}{}
\subsubsection*{Complex Dynamics of Multiphase Gas and Stars in X-ray Cooling Cores}
\addcontentsline{toc}{subsubsection}{Complex Dynamics of Multiphase Gas and Stars in X-ray Cooling Cores}
\textbf{Speaker:} Marie-Joëlle Gingras \\
\textbf{Affiliation:} University of Waterloo, Waterloo Centre for Astrophysics \textperiodcentered{} Student (BSc/MSc/PhD)\\
Studies of galaxy cluster cores have shown that radio-mechanical AGN feedback is needed to regulate the growth of massive galaxies. While the past decades have cemented the importance of AGN feedback in offsetting large-scale cooling, observations of multiphase gas and young stars associated with jets and X-ray cavities revealed that it can also locally promote cooling. However, our understanding of mechanical AGN feedback is still incomplete. I present optical IFU spectroscopy of the centres of four strong cooling cores. Few studies map stellar dynamics and resolve star formation histories alongside detailed nebular kinematics. This revealed peculiar velocities of the central galaxies relative to their nebular gas, later also observed in hot atmospheres by XRISM. Combined with X-ray and radio data, we compare the kinematics of different gas phases and stars. Abell 1835 and PKS 0745-191, both with powerful AGN feedback, contain two kinematically distinct phases of nebular gas: a quiescent phase and one churned-up by the AGN. Outflows aligned with X-ray cavities enabled the first direct measurements of cavity velocities from nebular gas. The low speeds suggest mass-loading from uplifted gas. The stellar kinematics and star formation histories of these systems further support AGN-stimulated cooling and star formation, from the extended young stellar populations in their cores and the detection in one system of young stars consistent with having cooled out of a co-spatial gas outflow. These results provide a multiphase view of radio-mechanical AGN feedback regulating cooling while shaping the kinematics and star formation histories of cluster cores.

\TopicNeedSpace{10\baselineskip}
\hypertarget{abs-53}{}
\subsubsection*{From Dwarfs to Giants: Galaxy and Group-Halo Mass Dependence of Black Hole Fueling \& Feedback at z\textasciitilde{}0}
\addcontentsline{toc}{subsubsection}{From Dwarfs to Giants: Galaxy and Group-Halo Mass Dependence of Black Hole Fueling \& Feedback at z\textasciitilde{}0}
\textbf{Speaker:} Sheila Kannappan \\
\textbf{Affiliation:} University of North Carolina at Chapel Hill \textperiodcentered{} Faculty (tenured)\\
I will present evidence for clear shifts in AGN fueling and feedback at key galaxy and group-halo mass scales, based on the uniquely comprehensive AGN census for the volume-limited z\textasciitilde{}0 RESOLVE and ECO surveys. This census combines traditional and new methods of optical and mid-IR detection, notably including both SF-AGN (low-metallicity/highly-star-forming AGN found using VO diagrams but not traditional BPT criteria) and Bonus AGN (star-forming AGN found using revised BPT criteria). Characteristic AGN types shift with galaxy and group-halo mass following correlations with star formation, with SF-AGN at the high-SF end (typical of dwarf, non-satellite galaxies), then Bonus AGN, then Composites, then Conventional AGN at the low-SF end (typical of massive or satellite galaxies). The volume-limited census design of RESOLVE and ECO, like a simulation box, yields AGN frequencies that directly constrain active black hole occupation fractions. Dividing RESOLVE+ECO into dwarf, transitional, and giant galaxies at the threshold and bimodality scales (galaxy mass \textasciitilde{}10\textasciicircum{}9.5 and 10\textasciicircum{}10.5 Msun, for centrals equivalent to halo mass \textasciitilde{}10\textasciicircum{}11.4 and 10\textasciicircum{}12 Msun), we find that: (1) inclusion of SF-AGN and Bonus AGN yields a \textasciitilde{}5x higher dwarf AGN frequency than typically reported, but nonetheless (2) AGN detections jump by a factor of \textasciitilde{}5 going from dwarf to transitional galaxies. Structural analysis shows that the sharp uptick in AGN detection at transitional masses is linked to "green nuggets": dense, compaction-formed centrals actively quenching from blue to red over the halo mass range of shock heating from \textasciitilde{}10\textasciicircum{}11.4 to 10\textasciicircum{}12 Msun. Feedback from dwarf/transitional AGN may add to such quenching, as we find clear group-integrated HI gas depletion for AGN-hosting halos in the same halo mass range. At higher masses we find weaker evidence for any AGN imprint on group HI content.

\section*{Theme 4: First Black Holes, Seeds and High-redshift}
\addcontentsline{toc}{section}{Theme 4: First Black Holes, Seeds and High-redshift}

\TopicNeedSpace{10\baselineskip}
\hypertarget{abs-65}{}
\subsubsection*{Quasar Clustering and Proximity Zones Constrain Early Supermassive Black Hole Growth}
\addcontentsline{toc}{subsubsection}{Quasar Clustering and Proximity Zones Constrain Early Supermassive Black Hole Growth}
\textbf{Speaker:} Elia Pizzati \\
\textbf{Affiliation:} Harvard University \textperiodcentered{} Postdoc\\
Understanding how supermassive black holes (SMBHs) formed and grew to billion-solar-mass scales by redshift z>7 remains a central challenge in astrophysics. Although many theoretical pathways have been proposed, the lack of robust empirical constraints has long prevented meaningful discrimination among competing formation scenarios.
In recent years, however, such constraints have begun to emerge. Proximity zone sizes in high-redshift quasar spectra now place tight limits on the duration of quasar activity episodes, while clustering measurements constrain the host halo masses and integrated duty cycles of quasars. For the first time, these observables allow us to move beyond instantaneous quantities — such as black hole mass and luminosity — and instead probe the full accretion histories of early quasars.
In this talk, I will present a new framework that leverages this emerging constraining power. The model generates synthetic SMBH accretion histories and quasar light curves coupled self-consistently to merger trees extracted from a large N-body cosmological simulation. This approach links individual SMBH growth tracks — including both accretion and mergers — and quasar spectra to population-level statistics such as clustering, black hole mass distributions, and luminosity functions.
Confronting these predictions with observations provides strong discriminatory power among competing formation scenarios: we find that sustained, episodic phases of super-critical accretion are required to simultaneously produce the most massive black holes by z>7 and reproduce the observed properties of the earliest quasars.


\hypertarget{abs-day-thursday}{}
\daypalettebanner{Thursday}{thursdaybg}{thursdayaccent}{July 2, 2026}

\section*{Theme 3: Accretion, Jets and Multi-messenger Signatures}
\addcontentsline{toc}{section}{Theme 3: Accretion, Jets and Multi-messenger Signatures}

\TopicNeedSpace{10\baselineskip}
\hypertarget{abs-2}{}
\subsubsection*{AGN feedback across cosmic time: from cosmic noon to reionization with JWST}
\addcontentsline{toc}{subsubsection}{AGN feedback across cosmic time: from cosmic noon to reionization with JWST}
\textbf{Speaker:} Stefano Carniani \invitedbadge{}\\

The epoch at z\textasciitilde{}3–6 marks a critical phase in galaxy evolution, when rapid transformations precede the peak of cosmic star formation and AGN activity. A key mechanism regulating this evolution is AGN feedback, though its impact at early times remains uncertain. Using JWST/NIRSpec IFU observations, I will present a study of AGN-driven outflows at high redshift. At z\textasciitilde{}3–6, we analyze a large spatially resolved sample of AGNs, mapping rest-frame optical mission lines to identify fast outflows and constrain their properties. We investigate their connection with both host galaxy properties and AGN luminosity, assessing their potential to regulate star formation.
We extend this analysis to the epoch of reionization, presenting new observations of the most distant quasars. These reveal powerful outflows on both nuclear and galactic scales, with mass outflow rates comparable to star formation rates, indicating that efficient feedback is already in place at very early times.
Finally, we place these results in a broader evolutionary context by comparing outflow properties across redshift and luminosity, and by examining the evolution of the black hole–stellar mass relation, whose increasing scatter at high redshift may reflect the bursty growth of both black holes and their host galaxies. Together, these findings constrain the onset and role of AGN feedback in shaping the earliest massive galaxies.

\TopicNeedSpace{10\baselineskip}
\hypertarget{abs-7}{}
\subsubsection*{Modeling Supermassive Black Hole Growth and Feedback in Galaxy Evolution}
\addcontentsline{toc}{subsubsection}{Modeling Supermassive Black Hole Growth and Feedback in Galaxy Evolution}
\textbf{Speaker:} Daniel Angles-Alcazar \invitedbadge{}\\

While observations across the electromagnetic spectrum are providing an increasingly detailed view of supermassive black holes and galaxies across cosmic time, theoretical models face the challenge of capturing the multi-scale physical processes that connect black holes to their host galaxies and cosmological environment. In this talk, I will review our current understanding of active galactic nucleus (AGN) fueling and feedback from cosmological hydrodynamic simulations, which provide the most direct framework for modeling supermassive black holes in a galaxy evolution context but rely on uncertain sub-grid prescriptions to represent unresolved physical processes. I will discuss (1) recent progress and ongoing challenges in modeling black hole seeding, dynamics, accretion, and feedback; (2) novel numerical hyper-refinement techniques that enable explicit modeling of AGN fueling and feedback from sub-parsec to circumgalactic medium scales; and (3) prospects for constraining theoretical models using observations from the nearby universe to the high-redshift frontier.

\TopicNeedSpace{10\baselineskip}
\hypertarget{abs-29}{}
\subsubsection*{Forged by Feedback: How AGN Feedback Models Shape Brightest Group Galaxies in Cosmological Simulations}
\addcontentsline{toc}{subsubsection}{Forged by Feedback: How AGN Feedback Models Shape Brightest Group Galaxies in Cosmological Simulations}
\textbf{Speaker:} Ruxin Barré \\
\textbf{Affiliation:} University of Victoria \textperiodcentered{} Student (BSc/MSc/PhD)\\
While most modern galaxy formation models can successfully reproduce global galaxy population statistics, many continue to struggle with capturing the properties and evolution of the most massive galaxies. These discrepancies are increasingly understood to arise from the modelling of active galactic nucleus (AGN) feedback and its critical role in regulating gas cooling and star formation. It therefore remains crucial to assess and refine the physical realism of AGN feedback prescriptions in cosmological simulations. Brightest group galaxies (BGGs), residing at the centres of galaxy groups, are shaped by supermassive black hole (SMBH)-driven cycles of heating and cooling, which provides an ideal laboratory for testing AGN feedback models. We use BGGs as probes of feedback physics, investigating their stellar properties across multiple cosmological simulations with distinct implementations of AGN feedback: purely thermal feedback in ROMULUS, differing calibrations of two-regime kinetic feedback in SIMBA and SIMBA-C, and the sophisticated, three-regime OBSIDIAN feedback model, which directly maps accretion flow geometry to the properties of SMBH-driven outflows. Our results reveal that the global properties and inferred quenching pathways of the simulated BGGs are highly sensitive to the strength and mechanism of their respective AGN feedback model. We contrast these predictions against observed BGGs in X-ray-selected COSMOS groups and find that the OBSIDIAN model achieves the highest agreement with observations, reproducing stellar property distributions, quenched fractions, and the evolution of star formation in increasingly massive systems. We find evidence that the OBSIDIAN and COSMOS BGG populations quench via a gradual decline in star formation with increasing stellar mass. This contrasts SIMBA and SIMBA-C, which exhibit rapid quenching linked to the onset of powerful jet feedback, and ROMULUS, whose thermal feedback cannot sufficiently suppress star formation. These results provide critical constraints for refining models of AGN feedback, with the success of OBSIDIAN signalling the start of a new generation of simulations with increasingly realistic depictions of SMBHs and their impact on galaxy evolution.

\TopicNeedSpace{10\baselineskip}
\hypertarget{abs-30}{}
\subsubsection*{Galaxies and Their Black Holes: Tracing and Harnessing a Cosmic Partnership}
\addcontentsline{toc}{subsubsection}{Galaxies and Their Black Holes: Tracing and Harnessing a Cosmic Partnership}
\textbf{Speaker:} Niel Brandt \\
\textbf{Affiliation:} Penn State University \textperiodcentered{} Faculty (tenured)\\
Supermassive black holes (SMBHs) and galaxies grow together over cosmic time, but the physical drivers of this joint evolution and SMBH feeding remain an open question. In this talk, I will summarize a statistical, population-based approach that links long-term SMBH growth to measurable galaxy properties. The key observable we use is the sample-averaged SMBH accretion rate (BHAR), inferred from sensitive X-ray surveys and interpreted as the long-term growth of SMBH populations. We compare BHAR against well-studied galactic quantities - including stellar mass, star-formation rate, and structural compactness - using large, well-controlled samples (millions of galaxies and thousands of active galactic nuclei) and partial-correlation methods that isolate the dominant connections.

Some broad conclusions are the following: (1) across z = 0-4 the generally strongest predictor of SMBH feeding and growth is galaxy stellar mass; (2) in systems dominated by a bulge component, the SMBH accretion rate tracks star formation closely, indicating synchronized SMBH and bulge growth; and (3) among star-forming galaxies, more compact systems show enhanced BHAR, plausibly because higher central gas densities feed accretion more efficiently. I will also describe how these empirical correlations can be folded into cosmological simulations to map SMBH mass assembly - via both accretion and mergers - and to explore consequences such as the evolving SMBH mass function, SMBH-galaxy scaling relations, feedback effects, and the drivers of the strong decline in SMBH accretion at z < 2.

\TopicNeedSpace{10\baselineskip}
\hypertarget{abs-31}{}
\subsubsection*{Mapping the Surrounding Ionized Gas of a Windy z > 6 Quasar with JWST/NIRSpec IFU}
\addcontentsline{toc}{subsubsection}{Mapping the Surrounding Ionized Gas of a Windy z > 6 Quasar with JWST/NIRSpec IFU}
\textbf{Speaker:} Hyunseop Choi \\
\textbf{Affiliation:} University of Michigan \textperiodcentered{} Postdoc\\
Quasars at z > 6 offer a unique opportunity to study supermassive black hole (SMBH) growth and feedback during the first billion years of cosmic history. I present new JWST/NIRSpec IFU observations of z \textasciitilde{} 6 quasars, focusing on J0923+0402 (z \textasciitilde{} 6.6), the highest-redshift BAL quasar with robust constraints on its outflow properties and the first laboratory where multiscale and multiphase quasar winds have been characterized through combined near-IR and sub-mm observations. This object hosts a powerful, warm-ionized BAL outflow near the center, with a much more extended cold-neutral outflow extending to kiloparsec scales.

Recent NIRSpec/IFU data provided the first spatially resolved view of rest-frame optical emission lines (e.g., [OIII], Balmer) in this windy quasar at cosmic dawn. We observed spatially extended warm ionized gas across kiloparsec scales and complex kinematics that directly connect the nuclear BAL wind to the host galaxy environment. When interpreted together with existing ALMA [C II] observations tracing neutral outflows and an extended halo, the JWST data revealed a complex, multiscale, and multiphase interaction between quasar-driven winds and the surrounding interstellar and circumgalactic media, only about \textasciitilde{}1 Gyr after the Big Bang.

I will discuss the coupling between compact BAL winds and galaxy-scale gas, the potential origin of extended emission-line regions, and the role of early SMBH feedback and inflow onto the SMBH in shaping the surrounding environment. I will also briefly present results from additional NIRSpec/IFU observations of z > 6 quasars, including one located in a galaxy overdensity and another showing strong BAL-driven winds.

\TopicNeedSpace{10\baselineskip}
\hypertarget{abs-32}{}
\subsubsection*{Uncovering Radio-Bright Tidal Disruption Events with ASKAP and the VLA}
\addcontentsline{toc}{subsubsection}{Uncovering Radio-Bright Tidal Disruption Events with ASKAP and the VLA}
\textbf{Speaker:} Hannah Dykaar \\
\textbf{Affiliation:} McGill University \textperiodcentered{} Postdoc\\
Tidal disruption events (TDEs) occur when a star ventures sufficiently close to a supermassive black hole, causing tidal forces to overcome the star's self-gravity. Radio observations can probe the size and velocity of these outflows as well as the density of the surrounding environment. Historically identified through shorter-wavelength observations, TDEs have benefited from recent advances in radio interferometry. In particular, the Australian Square Kilometre Array Pathfinder (ASKAP) and the Very Large Array (VLA), have ushered in a new era for TDE detection and follow-up. In this talk, I will present results from the first untargeted radio search for TDEs conducted with the recently completed Variables and Slow Transients (VAST) Pilot Survey from ASKAP. Unlike previous studies that relied on optical or X-ray discovery, our search identifies TDE candidates directly from radio variability and nuclear localization. We identify twelve radio-bright candidate TDEs, implying a volumetric rate of \textasciitilde{}0.8 Gpc⁻³ yr⁻¹, consistent with previous empirical estimates. This is the first constraint on the TDE rate derived from an untargeted radio survey. I also present a multiwavelength study of AT2022dbl, a repeating partial TDE, in which the stellar core survives the initial disruption and is subsequently disrupted again. Radio follow-up with the VLA places constraints on the density of the environment surrounding the central black hole. This analysis allows for the exploration of potential emission mechanisms and situates AT2022dbl within the broader context of TDEs, enhancing our understanding of partial TDEs and properties of their associated outflows. Together, these results demonstrate how wide-field radio surveys and detailed multiwavelength follow-up can jointly advance our understanding of TDE demographics, energetics, and outflow physics.

\TopicNeedSpace{10\baselineskip}
\hypertarget{abs-37}{}
\subsubsection*{Fueling, Winds, and Self-Regulation: SMBH Growth and Feedback from Cosmic Noon to the Local Universe with the FIRE simulations}
\addcontentsline{toc}{subsubsection}{Fueling, Winds, and Self-Regulation: SMBH Growth and Feedback from Cosmic Noon to the Local Universe with the FIRE simulations}
\textbf{Speaker:} Jonathan Mercedes Feliz \\
\textbf{Affiliation:} University of Connecticut \textperiodcentered{} Postdoc\\
Feedback from active galactic nuclei (AGN), powered by gas accretion onto supermassive black holes (SMBHs), is widely invoked as a key regulator of star formation in massive galaxies. However, the coupling between sub-kpc fueling, time-variable accretion, and multiphase outflows is hard to isolate observationally and difficult to capture in a single theoretical framework. High-resolution FIRE cosmological hydrodynamic simulations with a multiphase interstellar medium and explicit SMBH accretion and feedback models make it possible to follow these connections in time, from the gas supply in the circumgalactic medium (CGM) to nuclear inflow and the launching and impact of winds. In this presentation, I will summarize recent results focusing on three regimes that highlight the diversity of fueling and feedback across mass scale and cosmic time, from massive galaxies at their peak activity to Seyfert galaxies in the local universe. (1) I will first present a new nuclear fueling channel driving luminous quasars at cosmic noon based on the large pileup of gas in the inner CGM produced by a global galactic outflow that later reaccretes coherently onto the galaxy, resulting in significantly higher inflow rates than pre-outflow conditions that can directly supply the nuclear region and trigger a luminous quasar phase.  (2) I will then show that the resulting powerful AGN-driven winds can carve central cavities and suppress star formation in massive galaxies at cosmic noon primarily by preventing gas accretion into the ISM rather than by ejecting the existing star-forming gas reservoir. At the same time, strong AGN wind–ISM interactions can locally compress gas and result in positive feedback, including the rapid formation of ultra-dense, off-center stellar clumps. (3) Finally, I will discuss fueling and feedback self-regulation in simulations of Seyfert-like galaxies at low redshift, where we identify recurrent cycles linking increased inflow toward the accretion disk, enhanced SMBH accretion, feedback-driven clearing of the nuclear region, and subsequent re-fueling on \textasciitilde{}10–100 Myr timescales. Comparisons to the observed connection between AGN luminosity and nuclear gas deficit provide direct constraints on the characteristic spatial and temporal scales of SMBH self-regulation.

\TopicNeedSpace{10\baselineskip}
\hypertarget{abs-43}{}
\subsubsection*{From jet to filaments: a multi-scale and multi-wavelength view of AGN activity in M87}
\addcontentsline{toc}{subsubsection}{From jet to filaments: a multi-scale and multi-wavelength view of AGN activity in M87}
\textbf{Speaker:} Camille Poitras \\
\textbf{Affiliation:} Université Laval \textperiodcentered{} Student (BSc/MSc/PhD)\\
Active galactic nuclei (AGN) can launch relativistic jets that inject and redistribute energy in their surroundings. As the nearest brightest cluster galaxy in a cool-core cluster, M87 provides a unique laboratory to trace these processes in detail. In this talk, I will first present new constraints on the expansion of its well-known jet in X-rays with Chandra HRC-I, where emission from short-lived synchrotron electrons traces regions near ongoing particle acceleration. The sub-arcsecond imaging capability of HRC-I, combined with a 13-year baseline, enables measurements of proper motions and flux variability of individual knots, constraining their velocities, cooling behavior, and temporal evolution in a multi-wavelength context. On larger scales, jet-driven energy influences the extended warm filamentary nebula in M87. Using new integral-field spectroscopy observations from MEGARA (GTC) and SITELLE (CFHT), we provide the first complete views of its kinematics and excitation across the entire network. The gas exhibits turbulent motions and AGN-dominated excitation, and shows strong coupling to both the cold molecular phase and the cooler component of the hot intracluster medium. Together, these results provide new insight into AGN feedback across spatial scales, wavelengths, and timescales.

\TopicNeedSpace{10\baselineskip}
\hypertarget{abs-44}{}
\subsubsection*{AGN Feedback Signatures in Cosmological Simulations}
\addcontentsline{toc}{subsubsection}{AGN Feedback Signatures in Cosmological Simulations}
\textbf{Speaker:} Marine Prunier \\
\textbf{Affiliation:} UdeM \textperiodcentered{} Student (BSc/MSc/PhD)\\
AGN feedback plays a central role in shaping the thermodynamic structure of galaxy clusters, driving multiphase outflows that redistribute energy and metals across the intracluster medium (ICM). Capturing these complex processes in cosmological simulations is challenging, particularly the generation of observable small, kiloparsec-scale signatures such as X-ray cavities, shocks, and sound waves. In this talk, I will review recent advances in state-of-the-art cosmological simulations in reproducing AGN-driven structures that impact the hot ICM on kiloparsec scales. I will discuss methodologies for consistently comparing X-ray cavities and weak shocks in simulations with observations. Finally, I will highlight two key implications of these advancements. First, opening a new observational window to study the role of AGN feedback in a full cosmological context; and second, demonstrating that careful comparisons between these small-scale X-ray structures in simulations and observations are essential for validating SMBH subgrid models.

\TopicNeedSpace{10\baselineskip}
\hypertarget{abs-45}{}
\subsubsection*{Tidal Disruption Events: a probe of accretion onto supermassive black holes and their environments.}
\addcontentsline{toc}{subsubsection}{Tidal Disruption Events: a probe of accretion onto supermassive black holes and their environments.}
\textbf{Speaker:} Lauren Rhodes \\
\textbf{Affiliation:} McGill \textperiodcentered{} Postdoc\\
Tidal disruption events (TDEs) are rare, multi-wavelength transients that allow astronomers to probe otherwise quiescent supermassive black holes. Only a fraction of TDEs exhibit luminous radio emission, with fewer still appearing consistent with a relativistic jet. Despite their scarcity, the jetted TDEs discovered thus far have been rich sources of information, promoting the idea that future events, discovered by LSST and AXIS (if it launches), will be transformational to our understanding of jetted transients and supermassive black hole accretion. In this talk, I will present our comprehensive radio analysis of the jetted TDE AT2022cmc. I will demonstrate how our dense radio and sub-mm coverage allowed us to constrain both the large-scale geometry and microphysics of the jet. From there, I will use these results to motivate the importance of studying jetted TDEs as it provides a new path for studying the accretion mechanisms around supermassive black holes and a unique view of the inner parsec region of otherwise-quiescent galaxies.

\TopicNeedSpace{10\baselineskip}
\hypertarget{abs-46}{}
\subsubsection*{Bridging Scales: Coupling the galactic nucleus to the larger cosmic environment}
\addcontentsline{toc}{subsubsection}{Bridging Scales: Coupling the galactic nucleus to the larger cosmic environment}
\textbf{Speaker:} Kung-Yi Su Su \\
\textbf{Affiliation:} Northwestern University \textperiodcentered{} Research staff / Scientist\\
Recent advances in Multi-Zone GRMHD now allow us to bidirectionally track gas flows between galactic scales and black-hole event horizons within a single simulation. This enables, for the first time, a first-principles, GRMHD-informed AGN feedback model that can be directly implemented in galaxy simulations at resolvable scales. We test the model in an M87-like isolated-galaxy simulation—with boundary conditions matched to those used in the Multi-Zone GRMHD run—and in cosmological zoom-in simulations of a massive cluster of comparable mass. We infer the black-hole spin required to reproduce the accretion rate observed by the Event Horizon Telescope and outline the conditions under which this black-hole accretion/feedback mode dominates.

\TopicNeedSpace{10\baselineskip}
\hypertarget{abs-47}{}
\subsubsection*{Tracking supermassive black hole-galaxy coevolutionary histories in galaxy formation models}
\addcontentsline{toc}{subsubsection}{Tracking supermassive black hole-galaxy coevolutionary histories in galaxy formation models}
\textbf{Speaker:} Bryan Terrazas \\
\textbf{Affiliation:} Oberlin College \textperiodcentered{} Faculty (tenure-track )\\
Galaxy formation models rely on supermassive black holes (SMBHs) to suppress star formation and reproduce the observed census of massive galaxies. Since the amount of feedback released is directly dependent on the SMBH accretion rate history in these models, it is critical to accurately reproduce SMBH growth histories across the galaxy population. Current state-of-the-art models produce vastly different SMBH mass - stellar mass relations, indicating a need for observational constraints. I will present results from an empirical model of SMBH growth built on top of the UniverseMachine galaxy formation model that uses observationally-derived probability distributions for accretion rate as a function of galaxy properties. These results indicate a diversity of black hole growth histories: early, rapid growth for the most massive SMBHs in z = 0 quiescent galaxies, and late, gradual growth for lower mass SMBHs in z = 0 star-forming galaxies. We also find substantial evolution of the SMBH mass - stellar mass relation over time. I will also present results that merge our empirical model with a physical semi-analytic model, exploring how these accretion histories affect both mass and energy feedback from SMBHs.


\hypertarget{abs-day-friday}{}
\daypalettebanner{Friday}{fridaybg}{fridayaccent}{July 3, 2026}

\section*{Theme 2: Black Hole Growth and AGN Feedback}
\addcontentsline{toc}{section}{Theme 2: Black Hole Growth and AGN Feedback}

\TopicNeedSpace{10\baselineskip}
\hypertarget{abs-11}{}
\subsubsection*{A multi-messenger view of black holes in the event horizon-scale imaging era}
\addcontentsline{toc}{subsubsection}{A multi-messenger view of black holes in the event horizon-scale imaging era}
\textbf{Speaker:} Sera Markoff \invitedbadge{}\\
\textbf{Affiliation:} API/GRAPPA, University of Amsterdam and IoA/KICC, University of Cambridge\\
The last decade has witnessed spectacular developments in astrophysics and astroparticle physics, not least of which is the ability to directly image two supermassive black holes with global very long baseline interferometry (VLBI) project the Event Horizon Telescope (EHT). But how do these images connect to broader questions regarding the workings of accretion in general, or the launching of jets and particle acceleration, common to a wide variety of astrophysical phenomena? In this talk I will review how the combination of horizon-scale imaging with multi-wavelength and multi-messenger data, together with recent advances in theoretical modelling, have the potential to tackle some of the hairiest problems relating to black holes and compact objects in general. In particular I will discuss some recent results from the EHT multi-wavelength campaigns, as well as the theoretical interpretation and challenges so far. I will also try to give some perspective on what is (literally) on the horizon in the coming decade, as we move into an era of dynamical imaging and EHT-expansion projects, as well as science with the next-generation very-high-energy gamma-ray facility the Cherenkov Telescope Array Observatory and upgraded neutrino facilities, as well as GRAVITY+.

\TopicNeedSpace{10\baselineskip}
\hypertarget{abs-26}{}
\subsubsection*{Horizon-scale dynamics in the accretion flow of the supermassive black hole in Sagittarius A*}
\addcontentsline{toc}{subsubsection}{Horizon-scale dynamics in the accretion flow of the supermassive black hole in Sagittarius A*}
\textbf{Speaker:} Rohan Dahale \\
\textbf{Affiliation:} University of Toronto \textperiodcentered{} Postdoc\\
The Event Horizon Telescope (EHT) has imaged the time-averaged structure of the black hole at the Galactic Center, Sagittarius A* (Sgr A*), but resolving its temporal evolution enables fundamentally new science: directly observing how matter accretes onto the black hole and how plasma structures evolve on minute timescales near the event horizon. This provides tests of accretion physics impossible from static images alone. However, Sgr A*'s rapid variability on timescales of minutes, combined with sparse interferometric coverage, makes dynamic imaging extremely challenging.

We now aim to spatially resolve the time evolution of Sgr A* by reconstructing a minute-by-minute video of the black hole's immediate surroundings at horizon scale using April 11, 2017 EHT observations; the best dataset from the 2017 campaign, providing \textasciitilde{}3 hours of continuous coverage with ten simultaneous baselines. This day exhibited pronounced multi-wavelength variability: an X-ray flare (Chandra/NuSTAR) immediately before EHT observations, Q-U loops in ALMA polarimetric light curves, and quasi-periodic oscillations in EHT visibility-domain polarization, all suggesting coherent dynamics in the accretion flow.

To extract robust videos, we developed: (1) Multiple independent dynamic imaging algorithms using complementary approaches (neural fields, Bayesian inference, optimal transport regularization); (2) A comprehensive validation framework that rigorously tests whether these algorithms distinguish genuine astrophysical motion from reconstruction artifacts, using synthetic datasets ranging from simple geometric models to complex GRMHD simulations.

These validated reconstructions will reveal horizon-scale accretion dynamics and enable direct measurements of rotational plasma velocities and evolving brightness structures. Combined with multi-wavelength observations following the X-ray flare, this provides constraints on magnetic field geometry, plasma flow properties, and flare generation mechanisms. This work establishes a new observational frontier for understanding the physical processes driving variability in low-luminosity black hole accretion and testing theoretical predictions in the strong-gravity regime.

\TopicNeedSpace{10\baselineskip}
\hypertarget{abs-27}{}
\subsubsection*{The observations and science of the first black hole movie on event-horizon scales}
\addcontentsline{toc}{subsubsection}{The observations and science of the first black hole movie on event-horizon scales}
\textbf{Speaker:} Boris Georgiev \\
\textbf{Affiliation:} UArizona \textperiodcentered{} Faculty (tenured)\\
Several years of event-horizon images of a black hole have confirmed persistent strong gravity and a dynamic hot plasma that sources the jets which extend to galactic scales. However, models have a large parameter space and are routinely challenged by source variability, and interpretation is complicated by the entanglement of multiple astrophysical processes. Many open questions remain about the accretion flow, the origin of jet launching and acceleration, the magnetic field structure, the electron thermodynamics, and the origin of high-energy flares. In the spring of 2026, the Event Horizon Telescope will conduct a monitoring campaign of the black hole M87*. These observations will probe 10 to 360 gravitational timescales and are expected to show coherent motion of features on event-horizon scales. Temporal stacking will improve the dynamic range of the diffuse jet, and polarization will inform on the magnetic field. In this talk, I will describe these observations and detail the type of science they will enable. A velocity field promises an identification of the jet-launching region and the mechanism behind bright features, such as plasma hotspots, magnetized flux tubes, or wave-like perturbations. Similarly, magnetohydrodynamical models make strong predictions about the motion of rotating and launched features, allowing observations to more phenomenologically prefer certain mechanisms. Simultaneous multiwavelength observations could offer clues about electron acceleration in the rare chance of a concurrent high-energy flare.

\TopicNeedSpace{10\baselineskip}
\hypertarget{abs-28}{}
\subsubsection*{JWST’s First Detections of Mid-IR Flares from Sgr A*}
\addcontentsline{toc}{subsubsection}{JWST’s First Detections of Mid-IR Flares from Sgr A*}
\textbf{Speaker:} Daryl Haggard \\
\textbf{Affiliation:} McGill University/Trottier Space Institute \textperiodcentered{} Faculty (tenured)\\
Sgr A*, the black hole at the heart of the Milky Way Galaxy, has long been observed as a variable source. It displays high contrast X-ray flares and near-continuous variability at near IR wavelengths. Recently, we’ve detected the first flares at mid-IR wavelengths with the James Web Space Telescope’s MIRI IFU. I will present these new MIR flare detections and place them in the context of variability observed at other wavelengths, including by the EHT/SMA/ALMA (sub-mm), Chandra (X-ray), NuSTAR (hard X-ray) and others. Together these observations offer a deep, multi-wavelength, time-domain view of the closest supermassive black hole and allow us to constrain the physical and radiative processes that are driving variability near the event horizon. This JWST/MIRI program has also led to excellent Galactic Center science at larger scales and I’ll briefly describe improved solutions for Galactic extinction at MIR wavelengths, constraints on an IMBH in the Galactic Center, new abundance measurements for diffuse gas near Sgr A*, and new detections of a GC X-ray transient, MAXI J1744-294.

